\documentclass[10pt,a4paper]{amsart}
\usepackage{amsmath,amssymb,amsthm,graphicx}
\usepackage{verbatim}
\newenvironment{reference}{\comment}{\endcomment}
%\newenvironment{reference}{}{}
\newenvironment{slogan}{\comment}{\endcomment}
\newenvironment{history}{\comment}{\endcomment}
% For commutative diagrams we use Xy-pic
\usepackage[all]{xy}

% We use 2cell for 2-commutative diagrams.
\xyoption{2cell}
% Theorem environments.
%
\theoremstyle{plain}
\newtheorem{theorem}[subsection]{Theorem}
\newtheorem{proposition}[subsection]{Proposition}
\newtheorem{lemma}[subsection]{Lemma}

\theoremstyle{definition}
\newtheorem{definition}[subsection]{Definition}
\newtheorem{example}[subsection]{Example}
\newtheorem{exercise}[subsection]{Exercise}
\newtheorem{situation}[subsection]{Situation}

\theoremstyle{remark}
\newtheorem{remark}[subsection]{Remark}
\newtheorem{remarks}[subsection]{Remarks}

\numberwithin{equation}{subsection}

% Macros
%
\def\lim{\mathop{\mathrm{lim}}\nolimits}
\def\colim{\mathop{\mathrm{colim}}\nolimits}
\def\Spec{\mathop{\mathrm{Spec}}}
\def\Hom{\mathop{\mathrm{Hom}}\nolimits}
\def\Ext{\mathop{\mathrm{Ext}}\nolimits}
\def\SheafHom{\mathop{\mathcal{H}\!\mathit{om}}\nolimits}
\def\SheafExt{\mathop{\mathcal{E}\!\mathit{xt}}\nolimits}
\def\Sch{\mathit{Sch}}
\def\Mor{\mathop{\mathrm{Mor}}\nolimits}
\def\Ob{\mathop{\mathrm{Ob}}\nolimits}
\def\Sh{\mathop{\mathit{Sh}}\nolimits}
\def\NL{\mathop{N\!L}\nolimits}
\def\CH{\mathop{\mathrm{CH}}\nolimits}
\def\proetale{{pro\text{-}\acute{e}tale}}
\def\etale{{\acute{e}tale}}
\def\QCoh{\mathit{QCoh}}
\def\Ker{\mathop{\mathrm{Ker}}}
\def\Im{\mathop{\mathrm{Im}}}
\def\Coker{\mathop{\mathrm{Coker}}}
\def\Coim{\mathop{\mathrm{Coim}}}

% Boxtimes
%
\DeclareMathSymbol{\boxtimes}{\mathbin}{AMSa}{"02}

%
% Macros for moduli stacks/spaces
%
\def\QCohstack{\mathcal{QC}\!\mathit{oh}}
\def\Cohstack{\mathcal{C}\!\mathit{oh}}
\def\Spacesstack{\mathcal{S}\!\mathit{paces}}
\def\Quotfunctor{\mathrm{Quot}}
\def\Hilbfunctor{\mathrm{Hilb}}
\def\Curvesstack{\mathcal{C}\!\mathit{urves}}
\def\Polarizedstack{\mathcal{P}\!\mathit{olarized}}
\def\Complexesstack{\mathcal{C}\!\mathit{omplexes}}
% \Pic is the operator that assigns to X its picard group, usage \Pic(X)
% \Picardstack_{X/B} denotes the Picard stack of X over B
% \Picardfunctor_{X/B} denotes the Picard functor of X over B
\def\Pic{\mathop{\mathrm{Pic}}\nolimits}
\def\Picardstack{\mathcal{P}\!\mathit{ic}}
\def\Picardfunctor{\mathrm{Pic}}
\def\Deformationcategory{\mathcal{D}\!\mathit{ef}}
\usepackage{cleveref}
\crefname{originallemma}{original lemma}{original lemmas}
\Crefname{originallemma}{Original lemma}{Original lemmas}
\makeatletter
\newlabel{algebra-lemma-finite-extension-dim-1}{{02MJ}{source}{}{sourcepointer-1}{}}
\newlabel{algebra-lemma-finite-extension-dim-1@cref}{{[originallemma][0][]02MJ}{[1][0][]source}{}{}{}}
\makeatother
\crefname{originallemma}{original lemma}{original lemmas}
\Crefname{originallemma}{Original lemma}{Original lemmas}
\makeatletter
\newlabel{lemma-finite-in-codimension-one}{{02RM}{source}{}{sourcepointer-2}{}}
\newlabel{lemma-finite-in-codimension-one@cref}{{[originallemma][0][]02RM}{[1][0][]source}{}{}{}}
\makeatother
\crefname{originalreference}{original reference}{original references}
\Crefname{originalreference}{Original reference}{Original references}
\makeatletter
\newlabel{situation-setup}{{02QL}{source}{}{sourcepointer-3}{}}
\newlabel{situation-setup@cref}{{[originalreference][0][]02QL}{[1][0][]source}{}{}{}}
\makeatother
\title{Pushforward of a principal divisor along a dominant proper morphism equals the divisor of the norm}
\author{The Stacks Project Authors}
\date{stacks-project:chow.tex:local-snapshot-20261012; GNU FDL 1.2 or later}
\begin{document}
\maketitle
\noindent\textit{Editorial note.} All delivered mathematical text is the author's own, copied from the pinned Stacks Project chapter file chow.tex; no statement or proof step was rewritten, reordered or removed.\par\medskip
\noindent\textit{Editorial note.} Preamble provenance: the pinned chapter file opens with \textbackslash{}input\{preamble\} and carries no preamble text of its own, so the project's single Stacks-wide preamble staged at sources/stacks-chow/preamble.tex is excerpted verbatim as three spans, lines 8-12 (the verbatim package together with the project's own reference, slogan and history display-environment definitions, including the author's commented-out older form of the reference environment), lines 14-18 (Xy-pic setup) and lines 154-217 (the native theorem declarations and mathematical macros). The lines 8-12 span is required because a counted statement or its surrounding text may use those display environments, which the project typesets as comments. Nothing was added, removed or redefined; the wrapper class and the amsmath/amssymb/amsthm/graphicx base packages are supplied by the assembler.\par\medskip
\noindent\textit{Editorial note.} Same-chapter pointers are located at their \textbackslash{}label lines inside the pinned chapter file; cross-chapter pointers are located in the project's own tag list staged in the snapshot (sources/stacks-chow/tags.txt), at the verbatim lines carrying their canonical tags.\par\medskip
\noindent\textit{Allowed standard prerequisites (editorial).}\begin{itemize}
\item A dominant morphism of integral schemes of the same dimension is finite in codimension one. Same chapter chow.tex, label lemma-finite-in-codimension-one, invoked at proof line 2662. At the generic point of an integral closed subscheme of codimension one it lets the target be replaced by an affine neighbourhood over which the morphism is finite.
\item Norm compatibility of orders in a finite extension of dimension one. Algebra, label algebra-lemma-finite-extension-dim-1, invoked at proof line 2680. Supplies the equality between the order at a height one prime of the norm and the sum over the primes above it of the residue degrees times the orders, which is exactly the coefficient comparison.
\end{itemize}\medskip
\ifdefined\hypertarget\hypertarget{sourcepointer-1}{}\fi
\noindent\textit{Source reference.} Original lemma 02MJ: sources/stacks-chow/tags.txt, source lines 3446--3446. Navigation pointer for \textbackslash{}ref\{algebra-lemma-finite-extension-dim-1\} used inside the counted span (line 2680); the argument that consumes it is retained in the span.\par\medskip
\ifdefined\hypertarget\hypertarget{sourcepointer-2}{}\fi
\noindent\textit{Source reference.} Original lemma 02RM: sources/stacks-chow/pinned-original-source.tex, source lines 2524--2524. Navigation pointer for \textbackslash{}ref\{lemma-finite-in-codimension-one\} used inside the counted span (line 2662); the argument that consumes it is retained in the span.\par\medskip
\ifdefined\hypertarget\hypertarget{sourcepointer-3}{}\fi
\noindent\textit{Source reference.} Original reference 02QL: sources/stacks-chow/pinned-original-source.tex, source lines 1372--1372. Navigation pointer for \textbackslash{}ref\{situation-setup\} used inside the counted span (line 2646); the argument that consumes it is retained in the span.\par\medskip
\begin{lemma}
\label{lemma-proper-pushforward-alteration}
Let $(S, \delta)$ be as in Situation \ref{situation-setup}.
Let $X$, $Y$ be locally of finite type over $S$. Assume $X$, $Y$
are integral and $n = \dim_\delta(X) = \dim_\delta(Y)$.
Let $p : X \to Y$ be a dominant proper morphism.
Let $f \in R(X)^*$. Set
$$
g = \text{Nm}_{R(X)/R(Y)}(f).
$$
Then we have
$p_*\text{div}(f) = \text{div}(g)$.
\end{lemma}

\begin{proof}
Let $Z \subset Y$ be an integral closed subscheme of $\delta$-dimension
$n - 1$. We want to show that the coefficient of $[Z]$ in
$p_*\text{div}(f)$ and $\text{div}(g)$ are equal. We may apply
Lemma \ref{lemma-finite-in-codimension-one}
to the morphism $p : X \to Y$ and the generic point $\xi \in Z$.
Hence we may replace $Y$ by an
affine open neighbourhood of $\xi$ and assume that $p : X \to Y$ is finite.
Write $Y = \Spec(R)$ and $X = \Spec(A)$ with $p$ induced
by a finite homomorphism $R \to A$ of Noetherian domains which induces
an finite field extension $L/K$ of fraction fields.
Now we have $f \in L$, $g = \text{Nm}(f) \in K$,
and a prime $\mathfrak p \subset R$ with $\dim(R_{\mathfrak p}) = 1$.
The coefficient of $[Z]$ in $\text{div}_Y(g)$ is
$\text{ord}_{R_\mathfrak p}(g)$.
The coefficient of $[Z]$ in $p_*\text{div}_X(f)$ is
$$
\sum\nolimits_{\mathfrak q\text{ lying over }\mathfrak p}
[\kappa(\mathfrak q) : \kappa(\mathfrak p)]
\text{ord}_{A_{\mathfrak q}}(f)
$$
The desired equality therefore follows from
Algebra, Lemma \ref{algebra-lemma-finite-extension-dim-1}.
\end{proof}

\end{document}
